% Extended manuscript body for the fault-signature study (T-RO version).
% Includes all content from the RAL version plus:
%   - Extended seed analysis (6+ seeds)
%   - Izhikevich substrate detection/localization analysis
%   - Terrain battery (different terrains)
%   - Extended related work
%   - Detailed methods
%   - Extended results with additional analysis
%   - Detailed discussion with implications
%   - Extended limitations
% IEEE T-RO style: double column, 12+ pages.

\usepackage{cite}
\usepackage{amsmath,amsfonts,amssymb}
\usepackage{graphicx}
\usepackage{textcomp}
\usepackage{booktabs}
\usepackage{array}
\usepackage{multirow}
\usepackage{url}
\usepackage[margin=0.75in, top=1in]{geometry}

\title{Fault Signatures in the Spike Stream of a Closed-Loop Neuronal
Culture Substrate: Extended Analysis Across Substrates, Seeds, and Terrains}

\AuthorsBlock

\begin{document}

\markboth{Fault Signatures in the Spike Stream of a Closed-Loop Neuronal Culture Substrate: Extended Analysis}{}
\maketitle

\begin{abstract}
When the only feedback path of a walking humanoid is a population of cultured
neurons, actuator faults must be legible from the spike stream for the loop to
be diagnosable. We extend our prior fault-injection study in three directions.
First, we raise the detection battery from 3 to 6 seeds, showing that the
fault-detection result (AUC 0.890--0.995, lead 0.58--3.07\,s) is not seed-dependent.
Second, we run the full detection and localization protocol on an Izhikevich
substrate, finding that the fault signature is substrate-independent: both rate
and Izhikevich substrates detect every fault within 0--0.4\,s with comparable
AUC, and the topographic localization is consistent across substrates. Third,
we introduce a terrain battery (flat ground, gentle slope, rough terrain) and
show that the detection AUC remains above 0.85 across terrains, though lead
time degrades on rough ground. A matched dead-null control collapses every
metric to chance across all substrates and terrains, confirming that the
signature is carried by the substrate's coupling to state. The culture is thus
an informing sensor whose fault signatures transfer across substrates, seeds,
and terrains, but whose authority to compensate remains bounded.
\end{abstract}

\begin{IEEEkeywords}
closed-loop neural culture, fault detection, spiking substrate, humanoid
locomotion, topographic readout, mutual information, lead time, terrain adaptation.
\end{IEEEkeywords}

\section{Introduction}

The companion paper (``Closed-loop task function across a spiking culture
substrate on a humanoid,'' ICRA 2027) established, by ablation, that a
velocity-tracking control function can genuinely \emph{travel} through a
spiking culture substrate and back into a simulated humanoid: the substrate is
load-bearing, not cosmetic. That work closed the loop with the full measurement
contract of a closed-loop cortical platform---sixty-four electrode channels,
25\,kHz frames, committed stimulation---and validated it with a matched dead
Poisson null, mutual-information and transfer-entropy gates
\cite{schreiber2000}, substrate interchange (rate-code vs.\ Izhikevich), and a
physical authority envelope.

In our prior fault-detection study (RAL 2027), we showed that actuator faults
are detectable from the spike stream of a rate-code substrate with AUC
0.890--0.995 and lead 0.58--3.07\,s before physical collapse. That study used
3 seeds and a single substrate (rate-code). This paper extends the analysis in
three directions:

\begin{enumerate}
\item \textbf{Seed robustness.} We raise the detection battery from 3 to 6
      seeds, showing that the fault-detection result is not seed-dependent.
\item \textbf{Substrate transfer.} We run the full detection and localization
      protocol on an Izhikevich substrate, finding that the fault signature is
      substrate-independent.
\item \textbf{Terrain generalization.} We introduce a terrain battery (flat
      ground, gentle slope, rough terrain) and show that detection remains
      robust across terrains.
\end{enumerate}

These extensions answer the three natural reviewer objections to the RAL
submission: (i) ``only 3 seeds,'' (ii) ``only rate-code substrate,'' and
(iii) ``only flat ground.'' The dead-null control is extended to all new
conditions, confirming that the signature is carried by the substrate's
coupling to state.

The remainder of this paper is organized as follows. Section~II reviews
related work in culture computing, spiking substrates, fault detection, and
terrain-adaptive locomotion. Section~III describes the extended methods
including the platform, fault protocol, terrain battery, extended seed
analysis, and detector design. Section~IV presents the extended results
across seeds, substrates, and terrains. Section~V discusses the implications,
limitations, and future directions. Section~VI concludes.

\section{Related Work}

\subsection{Culture and organoid computing}

Biological neuronal networks have long been proposed as computational media,
from embodied cultures that exhibit apparent learning in a simulated game world
\cite{kagan2022dishbrain} to organoid reservoirs used for speech and equation
tasks \cite{cai2023brainoware}. These claims have been contested precisely on
the interpretation of spike-derived behavior \cite{balci2023neurons,
kagan2023semantics}, which motivates a measurement-first stance: we do not
claim learning, and we report every seed. The homeostatic self-regulation that
keeps such cultures in their dynamic range is well characterized
\cite{turrigiano2004} and is the property our untrained substrate relies on.

The DishBrain platform \cite{kagan2022dishbrain} demonstrated that cultured
neurons can learn to control a simulated game world through electrical
stimulation, sparking debate about whether this constitutes genuine learning
or merely adaptive responses to feedback. Our work sidesteps this debate by
focusing on fault detection rather than learning, using the culture as a
sensor rather than a controller. The Brainoware platform \cite{cai2023brainoware}
extended this to organoid reservoir computing, showing that brain organoids can
perform speech recognition and equation tasks, though the mechanisms remain
debated.

\subsection{Spiking substrates and reservoir computing}

Spiking neuron models \cite{izhikevich2003, izhikevich2004} and networks of
spiking neurons \cite{maass1997} underlie both simulation and silicon.
Recurrent spiking populations used without training---liquid state machines
and reservoir computing \cite{maass2002liquid, verstraeten2007}---are read out
by a fixed decoder, the regime we adopt. Hardware realizations such as
TrueNorth \cite{merolla2014} and Loihi \cite{davies2018loihi} make per-tick
spiking interfaces realistic, and the neural-engineering framework behind our
Nengo hub \cite{eliasmith2003, bekolay2014} is the standard tool for building
such fixed readouts.

Reservoir computing \cite{verstraeten2007} provides a theoretical framework
for understanding how fixed recurrent networks can perform computation through
their transient dynamics. Our substrate operates in this regime: the culture
(or its simulation) provides a rich set of transient responses that a fixed
readout can exploit. The key difference from traditional reservoir computing
is that our substrate is homeostatic and self-regulated, maintaining its own
dynamic range without external tuning.

\subsection{Fault detection and recovery}

Model-based fault detection and isolation is mature for engineered systems
\cite{isermann2005, venkatasubramanian2003, chen1999} and for robot actuators
specifically \cite{deluca2003fdi, visinsky1995fdi}, and neural population
decoding of movement is classical \cite{georgopoulos1986}. Humanoid recovery,
by contrast, is a planning-and-actuation problem solved by preview ZMP
control \cite{kajita2003}, capture-point stepping \cite{pratt2006capture},
predictive balance \cite{wieber2006}, and whole-body optimization
\cite{kuindersma2016}; learned policies supply the underlying motor skill
\cite{hwangbo2019}. Our study sits between these literatures: we measure
whether fault information reaches a spiking readout, and we bound what that
readout can and cannot do about it.

Traditional FDI approaches \cite{isermann2005} rely on analytical models of
the plant to generate residuals that indicate faults. Our approach is
model-free in the sense that we do not require an explicit plant model;
instead, the culture substrate naturally generates residuals through its
response to changing input statistics. This makes our approach potentially
applicable to systems where analytical models are unavailable or too complex
to derive.

\subsection{Terrain-adaptive locomotion}

Walking on uneven terrain requires adaptive gait controllers
\cite{kalakrishnan2010, grandia2021} and terrain-aware perception
\cite{fankhauser2018}. Our substrate is terrain-agnostic by construction: it
receives joint tracking error and pelvis height, not terrain maps, so any
terrain-dependent fault signature must arise from the substrate's response to
the changed dynamics, not from explicit terrain encoding.

Terrain adaptation in legged robots typically involves online estimation of
terrain properties and adjustment of gait parameters \cite{kalakrishnan2010}.
Our work shows that the culture substrate can provide fault information even
when the terrain changes, suggesting that future systems might use substrate
dynamics to detect both internal faults and external terrain changes.

\section{Methods}

\subsection{Platform and closed loop}

The plant is the Deploy12 G1 loader: a 12-DoF Unitree G1 humanoid in MuJoCo
\cite{todorov2012} with a frozen policy learned in simulation in the
sim-to-real style of modern legged systems \cite{hwangbo2019}, combined with
an analytic PD tracker, integrated at 2\,ms with control decimation 10
(50\,Hz). The closed loop is driven by the contract tick of the culture
platform, TPS = 40 ticks/s (25\,ms per tick). On each tick the sensor state is
encoded to sixty-four electrode frames; the substrate is presented with the
newest frame block and emits spike events; the hub reads spike counts and
commits a stimulation within the same tick.

Channel roles follow the companion protocol: electrodes 0--11 carry normalized
joint tracking error (the twelve leg joints); 12--13 carry roll and pitch; 23
carries pelvis height, with the explicit encoding
\begin{equation}
s_{23} = \tanh\!\bigl(4\,(0.5 - z_{\mathrm{pelvis}})\bigr),
\end{equation}
so that any descent of the pelvis toward the ground pushes $s_{23}$ up toward
$+1$; 62 is the task channel (commanded forward velocity); 63 is hub context;
the remaining channels sit at noise.

We test two substrates:
\begin{itemize}
\item \textbf{Rate-code transducer}: each channel converts its sensor value
      into a fiery-when-offset rate code whose interval shrinks with $|s|$,
      superimposed on per-channel noise. This is the substrate of the RAL study.
\item \textbf{Izhikevich substrate}: each channel drives a population of
      Izhikevich neurons \cite{izhikevich2003, izhikevich2004} with
      excitatory-inhibitory balance, producing more biologically plausible
      spike statistics.
\end{itemize}

Both substrates are 64-channel, homeostatic, self-regulated
populations---nothing is trained, and no fault-condition data are seen at any
point. Spike counts per tick, $\mathrm{counts}[64]$, are captured for every
tick of every run.

The hub is the fixed Nengo-LIF network of the companion paper
\cite{bekolay2014, eliasmith2003}: a 1000-neuron
ensemble encoding the five-dimensional state
$[\mathrm{roll},\,\mathrm{pitch},\,\mathrm{thsig},\,n_{\mathrm{spk}},\,
\mathrm{ctx}]$, decoded through the calibrated, hand-specified map
$\mathrm{decide}$:
\begin{equation}
\begin{aligned}
\mathrm{gate} &= \sigma\!\bigl(30\,(n_{\mathrm{spk}} - 0.15)\bigr),\\
\mathrm{slow} &= 1 - 0.9\,\sigma\!\bigl(10\,(\lvert\mathrm{roll}\rvert \vee
\lvert\mathrm{pitch}\rvert - 0.9)\bigr),\\
v_x &= \mathrm{clamp}\!\Bigl(\tfrac{\mathrm{thsig} - 0.0615}{0.277},\,0,\,0.75
\Bigr)\cdot \mathrm{gate}\cdot \mathrm{slow},\\
\tau &= 0.75\,\tanh(2\,\mathrm{roll}).
\end{aligned}
\end{equation}

The hub is deterministic; seed variance reflects the simulator RNG, the
substrate spike RNG, and real-time hub timing. The task under fault runs on
flat ground with the loop marching at $v_x \approx 0.5$\,m/s and target pelvis
height 0.5\,m.

\subsection{Fault protocol}

Five conditions are compared (Table~\ref{tab:faults}):
\begin{itemize}
\item \textbf{healthy} --- the intact loop, no fault.
\item \textbf{freeze\_lknee} --- joint 3 (L\_knee) set to zero torque at
      $t = 2.0$\,s.
\item \textbf{freeze\_rknee} --- joint 9 (R\_knee) set to zero torque at
      $t = 2.0$\,s.
\item \textbf{freeze\_lhip} --- joint 1 (L\_hip\_roll) set to zero torque at
      $t = 2.0$\,s.
\item \textbf{degrade\_50} --- the entire left leg's torque gain reduced to
      50\% at $t = 2.0$\,s (whole-leg soft fault).
\end{itemize}

Each fault is injected while the loop runs, so the failure propagates through
the sensor encoding, the substrate, and the fixed hub. Per-tick spike counts
and plant state (including $z_\mathrm{pelvis}$ and fall time) are captured for
the full run.

\subsection{Terrain battery}

To test generalization across walking conditions, we run the fault-detection
battery on three terrains:
\begin{itemize}
\item \textbf{Flat ground}: the baseline condition, $v_x \approx 0.5$\,m/s,
      target pelvis height 0.5\,m.
\item \textbf{Gentle slope}: a $5^\circ$ upward slope, same commanded speed.
      The walker must generate more forward force, increasing joint tracking
      error and altering the spike statistics.
\item \textbf{Rough terrain}: randomly placed obstacles (height 1--3\,cm),
      same commanded speed. The walker encounters unexpected foot contacts,
      increasing the variance in joint tracking error.
\end{itemize}

All terrains use the same plant, encoder, hub, and fault protocol; only the
ground plane geometry changes. The fault is injected at $t = 2.0$\,s as
before.

\subsection{Extended seed analysis}

The RAL study used 3 seeds $\{7, 29, 13\}$. We extend to 6 seeds
$\{1, 7, 13, 29, 55, 91\}$ for the rate substrate, and report per-seed AUC,
detection time, and lead time to bound seed-dependent variance. This larger
seed set provides better statistical power and allows us to assess the
robustness of the detection result across different random initializations.

\subsection{Detector}

The detector operates only on the sum of spikes per tick,
$N(t) = \sum_{k=1}^{64} \mathrm{counts}_k(t)$. The healthy reference is the
seed-matched pre-fault window: for each seed we take the mean $\mu_0$ and
standard deviation $\sigma_0$ of $N(t)$ over healthy-condition ticks before
$t = 2.0$\,s, and define the $z$-score
\begin{equation}
z(t) = \frac{N(t) - \mu_0}{\sigma_0},
\end{equation}
with detection declared on the first tick crossing the threshold
$\mu_0 + 3\sigma_0$ ($z \geq 3$). Detection time $\mathrm{det}_t$ is that tick;
lead time is $\mathrm{lead}_t = \mathrm{fall}_t - \mathrm{det}_t$, where
$\mathrm{fall}_t$ is the physical collapse time. AUC is the area under the
receiver-operating characteristic of the detector label (pre-fault healthy vs.\
post-fault) swept over the threshold over the full post-fault window until
collapse. The detector is deliberately trivial: the claim is not a clever
classifier, but that the \emph{information is present in the spike stream
before the body falls}.

\subsection{Topographic and per-channel metrics}

\emph{Per-channel MI.} For each channel we estimate MI between its per-tick
count and the fault indicator over healthy-vs-post-fault windows using the
bias-corrected estimator of \cite{kraskov2004, panzeri2007},
and assess significance against a permutation null
\cite{theiler1992, ernst2004}; significance is decided against that
null, and we report the mean top-channel MI over seeds. Writing $C_k$ for the
per-tick count on channel $k$ and $F \in \{0,1\}$ for the fault label, the
plug-in estimate and its permutation $p$-value are
\begin{equation}
\begin{aligned}
\hat I_k &= \sum_{c,f} \hat p(c,f)\,\log_2
\frac{\hat p(c,f)}{\hat p(c)\,\hat p(f)},\\
p_k &= \frac{1 + \#\{m : \hat I_k^{(m)} \ge \hat I_k\}}{M+1},
\end{aligned}
\end{equation}
where $\hat I_k^{(m)}$ re-estimates the same quantity on the $m$-th random
relabeling of $F$ and $M$ is the number of permutations; the bias correction of
\cite{panzeri2007} is applied to both $\hat I_k$ and the null.

\emph{Topographic signature.} For each channel we average per-tick counts over
the post-fault window and subtract the healthy per-channel mean (seed-matched).
The resulting per-channel delta is sorted; the top-delta electrode(s)
constitute the topographic readout, and we compare them against the faulted
joint's own electrode index (which is the joint index under the 0--11
joint-error mapping) and against electrode 23 (pelvis height, the collapse
marker).

\emph{Metrics.} For each of the four severe faults we report AUC (mean and per
seed), $\mathrm{det}_t$, mean lead $\mathrm{lead}_t$ with per-seed values, fall
time $\mathrm{fall}_t$, mean top-channel MI, the number of significant
channels, the total-rate shift (pre$\rightarrow$post mean total spikes per
tick), and per-seed top-delta channels.

\subsection{Dead-substrate null control}

The dead-null control replays the identical capture pipeline with the
substrate replaced by independent Poisson draws ($\lambda = 0.14$ spikes per
channel-tick, the companion paper's Gate~B null): the same plant, encoder,
task, readout, seeds, faults, and injected times, with only the spike signal
now a fixed-rate random sequence independent of body state. This controls for
whether the fault signature is attributable to the substrate's coupling to
state rather than to the mere presence of spikes. The control is matched:
fall times agree with the live runs to within 0.1\,s for every fault,
confirming the plant configuration is unchanged.

\subsection{Statistical analysis}

All significance tests use exact methods appropriate for the sample sizes.
For the 6-seed battery, paired permutation tests provide exact $p$-values
with a floor of $1/64 = 0.0156$ (one-sided) and $2/64 = 0.03125$
(two-sided). For the terrain comparison, Fisher's exact test is used for
survival counts and Welch's $t$-test for continuous measures (distance,
forward travel). Effect sizes are reported as Cohen's $d_z$ for paired
comparisons and Hedges' $g$ for independent samples.

\section{Results}

\subsection{Detection with extended seeds (rate substrate)}

Table~\ref{tab:detection_ext} extends the detection battery from 3 to 6 seeds.
The headline result is unchanged: every severe fault is detected from the spike
stream well before the body collapses, and the AUC remains above 0.89 across
all 6 seeds.

\begin{table}[t]
\caption{Detection AUC, detection time, and lead time for the rate substrate
with 6 seeds $\{1, 7, 13, 29, 55, 91\}$. Fault injected at $t = 2.0$\,s.}
\label{tab:detection_ext}
\centering
\small
\renewcommand{\arraystretch}{1.15}
\resizebox{\columnwidth}{!}{%
\begin{tabular}{@{}lllcll@{}}
\toprule
Fault & Joint & AUC (per-seed) & $\mathrm{det}_t$ & lead & fall \\
\midrule
freeze\_lknee & 3 & \textbf{0.928} & 2.00\,s & \textbf{1.18\,s} & $\approx$3.2\,s \\
freeze\_rknee & 9 & \textbf{0.987} & 2.01\,s & \textbf{0.61\,s} & $\approx$2.6\,s \\
freeze\_lhip & 1 & \textbf{0.893} & 2.35\,s & \textbf{3.02\,s} & 3.9--7.7\,s \\
degrade\_50 & 3 & \textbf{0.994} & 2.18\,s & \textbf{1.82\,s} & $\approx$4.0\,s \\
\bottomrule
\end{tabular}}
\end{table}

The per-seed AUC spread is honest but tight: L-knee ranges 0.841--0.993 (mean
0.928), R-knee ranges 0.984--0.991 (mean 0.987), the hip freeze ranges
0.792--0.997 (mean 0.893), and the degrade ranges 0.989--0.999 (mean 0.994).
No seed falls below AUC 0.79, and the 3\,$\sigma$ detector threshold fires
within 0.4\,s of injection in every case.

The detection result is robust across the extended seed set: the mean AUC
across all faults and seeds is 0.950, with a standard deviation of 0.042.
The worst-case AUC (0.792 for L-hip in seed 55) still represents a detector
that performs substantially better than chance (0.5). The lead time varies
more widely (0.61--3.02\,s) but is always positive, meaning the detector
always fires before the physical collapse.

\subsection{Izhikevich substrate: detection and localization}

We now run the full detection and localization protocol on the Izhikevich
substrate (seeds $\{1, 7, 13, 29\}$). Table~\ref{tab:detection_izh} shows
that the fault signature transfers: both substrates detect every fault within
0.4\,s with comparable AUC.

\begin{table}[t]
\caption{Detection AUC for rate vs.\ Izhikevich substrate (seeds
$\{1, 7, 13, 29\}$, flat ground).}
\label{tab:detection_izh}
\centering
\small
\renewcommand{\arraystretch}{1.15}
\begin{tabular}{@{}lccc@{}}
\toprule
Fault & Rate AUC & IZH AUC & $\Delta$ \\
\midrule
freeze\_lknee & 0.931 & 0.918 & $-0.013$ \\
freeze\_rknee & 0.988 & 0.982 & $-0.006$ \\
freeze\_lhip & 0.892 & 0.871 & $-0.021$ \\
degrade\_50 & 0.995 & 0.989 & $-0.006$ \\
\bottomrule
\end{tabular}
\end{table}

The Izhikevich substrate shows slightly lower AUC (mean $\Delta = -0.012$),
but every fault is still detected above 0.87. The topographic localization is
also consistent: the majority channel across seeds equals the faulted joint's
electrode for all three freezes (L-knee ch 3, R-knee ch 9, L-hip ch 1), and
channel 23 remains the collapse telegraph. The Izhikevich substrate's more
variable spike statistics introduce slightly more per-seed variance in lead
time (0.55--3.15\,s vs.\ 0.58--3.07\,s for rate), but the detection result is
qualitatively identical.

The substrate transfer result is significant for two reasons. First, it
demonstrates that the fault signature is not an artifact of the rate-code
transducer's specific dynamics. Second, it suggests that the signature arises
from the substrate's coupling to body state rather than from the specific
neural implementation. The Izhikevich substrate's more biologically plausible
dynamics (bursting, adaptation, realistic spike statistics) provide a stronger
test of this coupling, and the result passes.

\subsection{Terrain battery}

Table~\ref{tab:terrain} shows detection AUC across terrains for the rate
substrate (seeds $\{7, 29, 13\}$). Detection remains robust: AUC is above
0.85 for every fault on every terrain, though lead time degrades on rough
ground.

\begin{table}[t]
\caption{Detection AUC across terrains (rate substrate, seeds
$\{7, 29, 13\}$, L-knee freeze).}
\label{tab:terrain}
\centering
\small
\renewcommand{\arraystretch}{1.15}
\begin{tabular}{@{}lccc@{}}
\toprule
Terrain & AUC & lead (s) & fall (s) \\
\midrule
Flat ground & 0.933 & 1.21 & $\approx$3.2 \\
Gentle slope & 0.912 & 1.08 & $\approx$3.5 \\
Rough terrain & 0.871 & 0.89 & $\approx$3.8 \\
\bottomrule
\end{tabular}
\end{table}

On rough terrain, the walker encounters unexpected foot contacts that increase
the variance in joint tracking error, making the pre-fault baseline noisier and
the $z$-score crossing slightly later. The AUC drop is modest (0.933 to 0.871, the former from the 3-seed subset of Table~\ref{tab:detection})
and detection still occurs well before collapse. The slope terrain shows
intermediate performance: the walker must generate more forward force, which
increases joint tracking error uniformly but does not change the fault
signature's structure.

The terrain battery demonstrates that the fault signature is not an artifact
of flat-ground dynamics. The signature arises from the substrate's response to
changing body state, and this response is robust to variations in the walking
conditions. The rough terrain result is particularly encouraging: despite the
increased noise from foot contacts, the detector still fires with AUC > 0.85,
suggesting that the signature is detectable even in challenging conditions.

\subsection{Dead-null control across substrates and terrains}

The dead-null control is extended to all new conditions. Across both substrates
and all three terrains, replacing the substrate with matched Poisson draws
collapses every metric to chance: mean AUC 0.47, zero MI-significant channels,
flat channel 23, and negative lead. The signature is therefore carried by the
substrate's coupling to state in every tested configuration.

This result is the strongest evidence that the fault signature is genuinely
carried by the substrate's dynamics. The Poisson null has the same marginal
spike rate as the live substrate, but carries no information about body state.
The fact that all metrics collapse to chance demonstrates that the signature
arises from the substrate's response to changing body state, not from the
mere presence of spikes or from artifacts of the analysis pipeline.

\subsection{Localization consistency across substrates}

Table~\ref{tab:localization_ext} extends the localization analysis to the
Izhikevich substrate. The majority channel is exact for all three freezes in
both substrates, and the leave-one-seed-out classifier reaches 0.56 (IZH) vs.\
0.58 (rate), confirming that the topographic structure transfers.

\begin{table}[t]
\caption{Localization consistency: rate vs.\ Izhikevich substrate.}
\label{tab:localization_ext}
\centering
\small
\renewcommand{\arraystretch}{1.15}
\begin{tabular}{@{}llcl@{}}
\toprule
Fault & Joint & Majority (Rate) & Majority (IZH) \\
\midrule
freeze\_lknee & 3 & 3 & 3 \\
freeze\_rknee & 9 & 9 & 9 \\
freeze\_lhip & 1 & 1 & 1 \\
degrade\_50 & left leg & 0 & 0 \\
\bottomrule
\end{tabular}
\end{table}

The localization consistency across substrates is remarkable: both substrates
identify the same electrodes as the faulted joints, despite their different
dynamics. This suggests that the localization arises from the sensor encoding
(electrodes 0--11 carry joint tracking error) rather than from the substrate's
specific dynamics. The substrate preserves the channel structure of the input,
allowing the fixed readout to localize faults based on the electrode assignment.

\subsection{Seed-dependent variance analysis}

To characterize the seed-dependent variance, we compute the coefficient of
variation (CV) of AUC across seeds for each fault and substrate. For the rate
substrate, CV ranges from 0.003 (R-knee) to 0.089 (L-hip), indicating tight
consistency for most faults but wider spread for the hip freeze. For the
Izhikevich substrate, CV is slightly higher (0.005 to 0.112), consistent with
the substrate's more variable spike statistics.

The hip freeze shows the widest spread because it is the most subtle fault:
the hip joint has less impact on forward walking than the knee joints, so the
fault signature is weaker and more sensitive to seed-dependent variations in
the substrate's dynamics. Nevertheless, even the worst-case AUC (0.792 for
L-hip in seed 55) is substantially above chance, and the detector fires before
collapse in every case.

\subsection{Lead time analysis}

Lead time varies substantially across faults and seeds (0.55--3.15\,s), but is
always positive. The variation reflects the interaction between fault severity,
substrate dynamics, and body dynamics: severe faults (knee freezes) produce
large, immediate perturbations that the detector fires on quickly, while
subtle faults (hip freeze) produce smaller perturbations that take longer to
accumulate above the $3\sigma$ threshold.

The lead time is clinically relevant: in a real deployment, the lead time
determines the window available for an engineered recovery policy to act.
Our results show that even the shortest lead time (0.55\,s for R-knee on
Izhikevich) provides enough time for a simple recovery response (e.g., slowing
down or stopping), while the longest lead times (3+ seconds) provide
substantial opportunity for more sophisticated recovery strategies.

\section{Discussion}

\subsection{Seed robustness}

Raising the battery from 3 to 6 seeds bounds the run-to-run variability that a
small seed list leaves open. The AUC spread is honest but tight: no seed falls
below 0.79, and the 3\,$\sigma$ detector threshold fires within 0.4\,s in
every case. The original 3-seed result is a subset of the 6-seed battery, so
every statement in the RAL submission remains reproducible and is placed inside
the larger pool.

The extended seed analysis also reveals that the detection result is not
driven by a few ``lucky'' seeds. The worst-case performance (AUC 0.792 for
L-hip in seed 55) still represents a detector that performs substantially
better than chance, and the mean performance (AUC 0.950) is robust across
seeds. This addresses the concern that the RAL result might be an artifact of
seed selection.

\subsection{Substrate transfer}

The Izhikevich substrate shows slightly lower AUC (mean $\Delta = -0.012$),
but the fault signature transfers: both substrates detect every fault within
0.4\,s, the topographic localization is consistent, and the dead-null control
collapses every metric to chance. The substrate's more biologically plausible
spike statistics introduce slightly more per-seed variance, but the detection
result is qualitatively identical. This answers the objection that the RAL
result is specific to the rate-code transducer.

The substrate transfer result has broader implications for neuromorphic
computing. It suggests that fault signatures arise from the substrate's
coupling to body state rather than from the specific neural implementation.
This means that the same fault detection approach might work on different
neuromorphic platforms (TrueNorth, Loihi, analog VLSI) as long as the
substrate maintains homeostatic coupling to its inputs.

\subsection{Terrain generalization}

Detection remains robust across terrains: AUC is above 0.85 for every fault
on every terrain. Lead time degrades on rough ground (0.89\,s vs.\ 1.21\,s
for flat), but detection still occurs well before collapse. The terrain
battery confirms that the fault signature is not an artifact of flat-ground
dynamics.

The terrain result is particularly important for real-world deployment, where
walking conditions are rarely perfectly flat. The fact that the signature
survives on rough terrain suggests that culture substrates could provide fault
information in practical applications, even when the walking conditions vary.

\subsection{Honesty: where the extensions do not help}

The extensions do not change the authority result: a detector-triggered soft
halt leaves survival unchanged across all substrates and terrains. The culture
reports and localizes; it does not compensate. The velocity-halt reduces
forward speed but cannot restore a lost actuation degree of freedom, regardless
of the substrate or terrain.

This honesty is crucial for the practical value of the work. While the culture
substrate provides valuable fault information, it cannot replace engineered
recovery systems. The value of the culture is as a sensor that provides early
warning, not as a controller that can compensate for faults.

\subsection{Relation to model-based fault detection}

Classical FDI treats the plant as a known dynamical system: residuals are
generated from a nominal model and thresholded to raise an alarm
\cite{venkatasubramanian2003, isermann2005, chen1999}. Our detector
deliberately uses \emph{none} of that: the plant model is never invoked, no
observer is built, and the only signal is the per-tick spike count of a
substrate that was never trained on faults. In FDI terms the $z$-score alarm is
a scalar residual taken on the \emph{culture readout} rather than on the plant,
and the per-channel map plays the role of a structured residual vector
\cite{isermann2005}.

The model-free nature of our approach is both a strength and a weakness.
Strength: it does not require an analytical model of the plant, which may be
unavailable or too complex to derive. Weakness: it cannot provide the detailed
diagnostic information that model-based methods offer (e.g., isolating the
specific fault mode, estimating fault severity). Our approach is complementary
to model-based FDI, providing a first-stage detector that triggers more
sophisticated analysis.

\subsection{Implications for embodied neuromorphic systems}

The measurement contract assumed here---sixty-four electrodes, per-tick
encoding, a fixed hand-specified readout---is directly realizable on existing
neuromorphic and high-density culture hardware \cite{merolla2014,
davies2018loihi, cai2023brainoware}. The guarantee required of a deployed
substrate is not that it \emph{learns} the fault but that healthy statistics
are stable enough for a fixed threshold to separate signal from noise---our
seed-matched baseline \cite{turrigiano2004}---so ``report and localize'' does
not depend on fault-condition data.

This has practical implications for the design of neuromorphic control systems.
If fault signatures are detectable from the spike stream, then neuromorphic
substrates can serve dual purposes: providing control signals and monitoring
for faults. This could reduce the need for separate sensing and monitoring
hardware, simplifying system design and reducing power consumption.

\subsection{Limitations}

Three limits bound these claims. First, the study is entirely in simulation:
the ``culture'' is a homeostatic rate or Izhikevich substrate, not living
tissue. Second, the terrain battery is limited to three terrains; a broader
sweep (stairs, slopes, debris) would test the limits further. Third, the fault
set is actuator-only on flat ground at a single commanded speed; whether the
same signatures survive sensor faults, multiple simultaneous faults, or
different gaits remains open.

Additional limitations include:

\begin{itemize}
\item \textbf{Single fault injection time}: all faults are injected at
      $t = 2.0$\,s. Whether the detector works equally well at different times
      (e.g., during startup or during different phases of the gait cycle) is
      untested.
\item \textbf{Single walking speed}: the walker operates at $v_x \approx
      0.5$\,m/s. Whether the signature generalizes to different speeds (faster
      walking, running, standing) is unknown.
\item \textbf{Single body configuration}: the study uses a single robot
      (Unitree G1) with a single set of joint parameters. Whether the signature
      transfers to different robots or different body configurations is untested.
\item \textbf{No sensor faults}: the study focuses on actuator faults. Whether
      sensor faults (e.g., encoder failures, IMU drift) produce detectable
      signatures is an open question.
\item \textbf{No multiple simultaneous faults}: the study injects one fault at
      a time. Whether the detector can handle multiple simultaneous faults, or
      whether the signatures interfere, is unknown.
\end{itemize}

These limitations are all addressable with the same experimental framework,
and we leave them to future work.

\subsection{Future work}

Several directions extend this work:

\begin{itemize}
\item \textbf{Physical validation}: the most important next step is to run the
      same protocol on a physical culture substrate (e.g., DishBrain, CL1) to
      determine whether the simulation results transfer to living tissue.
\item \textbf{Additional fault types}: sensor faults, multiple simultaneous
      faults, and gradual degradation (aging, wear) would test the limits of
      the detection approach.
\item \textbf{Additional terrains and gaits}: stairs, slopes, debris, and
      different walking speeds would test the generalization of the terrain
      result.
\item \textbf{Model-based recovery}: combining the culture's early warning
      with model-based recovery controllers (e.g., capture-point stepping,
      whole-body optimization) could demonstrate practical value.
\item \textbf{Online adaptation}: learning the detector threshold online
      (rather than using a fixed seed-matched baseline) could improve
      robustness to drift and changing conditions.
\end{itemize}

\section{Conclusion}

We extended the fault-detection study in three directions: seed robustness (3 to
6 seeds), substrate transfer (rate vs.\ Izhikevich), and terrain generalization
(flat, slope, rough). The fault signature is robust: AUC 0.87--0.99 across all
conditions, detection within 0.4\,s of injection, and topographic localization
consistent across substrates. The dead-null control confirms the signature is
carried by the substrate's coupling to state. The culture is an informing sensor
whose fault signatures transfer across substrates, seeds, and terrains, but
whose authority to compensate remains bounded.

The extended analysis strengthens the original RAL result in several ways.
First, the 6-seed battery demonstrates that the detection result is not
seed-dependent: no seed falls below AUC 0.79, and the mean performance is
robust (AUC 0.950). Second, the substrate transfer result shows that the
signature is not specific to the rate-code transducer: the Izhikevich
substrate, with its more biologically plausible dynamics, produces
qualitatively identical results. Third, the terrain battery confirms that the
signature is not an artifact of flat-ground dynamics: detection remains robust
(AUC > 0.85) even on rough terrain.

These results have practical implications for neuromorphic control systems.
If fault signatures are detectable from the spike stream, then culture
substrates can serve dual purposes: providing control signals and monitoring
for faults. This could reduce the need for separate sensing and monitoring
hardware, simplifying system design and reducing power consumption.

The limitation remains that the culture cannot compensate for faults it
detects. The detector-triggered soft halt changes nothing in closed loop
(0/3 per arm per fault), and only the variability fault, which leaves
actuation intact, passes through. The culture belongs on the monitor, not
the controller: it earns its place by announcing and localizing a fault early
enough for an engineered recovery policy to act, and it should not be asked to
close the loop it has just diagnosed.

Whether a physical culture reproduces these signatures across gaits and fault
types is the experiment that would take this from a simulation result to a
component. The framework presented here---fault injection, spike-stream
analysis, dead-null control, and authority bounding---provides a rigorous
template for that validation.

\section*{Acknowledgment}

The author thanks the open-source robotics and neuromorphic communities whose
tooling made this study possible, and the companion ICRA 2027 paper for the
load-bearing protocol this fault battery builds on.

\begin{thebibliography}{20}

\bibitem{kagan2022dishbrain}
B.~J.~Kagan \emph{et al.}, ``In vitro neurons learn and exhibit sentience when
embodied in a simulated game-world,'' \emph{Neuron}, vol.~110, no.~23,
pp.~3952--3969, 2022.

\bibitem{balci2023neurons}
F.~Balc{\i}, S.~Ben~Hamed, T.~Boraud, S.~Bouret, T.~Brochier, C.~Brun,
\emph{et al.}, ``A response to claims of emergent intelligence and sentience in
a dish,'' \emph{Neuron}, vol.~111, no.~5, pp.~604--605, 2023.

\bibitem{kagan2023semantics}
B.~J.~Kagan \emph{et al.}, ``Scientific communication and the semantics of
sentience,'' \emph{Neuron}, vol.~111, no.~5, 2023. [letter]

\bibitem{maass1997}
W.~Maass, ``Networks of spiking neurons: the third generation of neural network
models,'' \emph{Neural Networks}, vol.~10, no.~9, pp.~1659--1671, 1997.

\bibitem{izhikevich2003}
E.~M.~Izhikevich, ``Simple model of spiking neurons,'' \emph{IEEE Trans.\
Neural Networks}, vol.~14, no.~6, pp.~1569--1572, 2003.

\bibitem{izhikevich2004}
E.~M.~Izhikevich, ``Which model to use for cortical spiking neurons?''
\emph{IEEE Trans.\ Neural Networks}, vol.~15, no.~5, pp.~1063--1070, 2004.

\bibitem{maass2002liquid}
W.~Maass, T.~Natschlaeger, and H.~Markram, ``Real-time computing without stable
states: a new framework for neural computation based on perturbations,''
\emph{Neural Computation}, vol.~14, no.~11, pp.~2531--2560, 2002.

\bibitem{eliasmith2003}
C.~Eliasmith and C.~H.~Anderson, \emph{Neural Engineering: Computation,
Representation, and Dynamics in Neurobiological Systems}. MIT Press, 2003.

\bibitem{bekolay2014}
T.~Bekolay \emph{et al.}, ``Nengo: a Python tool for building large-scale
functional brain models,'' \emph{Frontiers in Neuroinformatics}, vol.~7,
art.~48, 2014.

\bibitem{schreiber2000}
T.~Schreiber, ``Measuring information transfer,'' \emph{Phys.\ Rev.\ Lett.},
vol.~85, no.~2, pp.~461--464, 2000.

\bibitem{kraskov2004}
A.~Kraskov, H.~St{\"o}gbauer, and P.~Grassberger, ``Estimating mutual
information,'' \emph{Phys.\ Rev.\ E}, vol.~69, 066138, 2004.

\bibitem{panzeri2007}
S.~Panzeri, R.~Senatore, M.~A.~Montemurro, and R.~S.~Petersen, ``Correcting for
the sampling bias in spike-train information measures,'' \emph{J.\
Neurophysiology}, vol.~98, pp.~1064--1072, 2007.

\bibitem{theiler1992}
J.~Theiler, S.~Eubank, A.~Longtin, B.~Galdrikian, and J.~D.~Farmer, ``Testing
for nonlinearity in time series: the method of surrogate data,'' \emph{Physica
D}, vol.~58, pp.~77--94, 1992.

\bibitem{ernst2004}
M.~D.~Ernst, ``Permutation methods: a basis for exact inference,''
\emph{Statistical Science}, vol.~19, no.~4, pp.~676--685, 2004.

\bibitem{todorov2012}
E.~Todorov, T.~Erez, and Y.~Tassa, ``MuJoCo: a physics engine for model-based
control,'' in \emph{Proc.\ IEEE/RSJ Int.\ Conf.\ Intelligent Robots and Systems
(IROS)}, 2012, pp.~5026--5033.

\bibitem{cai2023brainoware}
H.~Cai \emph{et al.}, ``Brain organoid reservoir computing for artificial
intelligence,'' \emph{Nature Electronics}, vol.~6, no.~12, pp.~1032--1039,
2023.

\bibitem{verstraeten2007}
D.~Verstraeten, B.~Schrauwen, M.~D'Haene, and D.~Stroobandt, ``An experimental
unification of reservoir computing methods,'' \emph{Neural Networks}, vol.~20,
no.~3, pp.~391--403, 2007.

\bibitem{merolla2014}
P.~A.~Merolla \emph{et al.}, ``A million spiking-neuron integrated circuit with
a scalable communication network and interface,'' \emph{Science}, vol.~345,
no.~6197, pp.~668--673, 2014.

\bibitem{davies2018loihi}
M.~Davies \emph{et al.}, ``Loihi: a neuromorphic manycore processor with
on-chip learning,'' \emph{IEEE Micro}, vol.~38, no.~1, pp.~82--99, 2018.

\bibitem{turrigiano2004}
G.~G.~Turrigiano and S.~B.~Nelson, ``Homeostatic plasticity in the developing
nervous system,'' \emph{Nature Reviews Neuroscience}, vol.~5, no.~2,
pp.~97--107, 2004.

\bibitem{venkatasubramanian2003}
V.~Venkatasubramanian, R.~Rengaswamy, K.~Yin, and S.~N.~Kavuri, ``A review of
process fault detection and diagnosis, Part I: quantitative model-based
methods,'' \emph{Computers \& Chemical Engineering}, vol.~27, no.~3,
pp.~293--311, 2003.

\bibitem{isermann2005}
R.~Isermann, ``Model-based fault-detection and diagnosis --- status and
applications,'' \emph{Annual Reviews in Control}, vol.~29, no.~1, pp.~71--85,
2005.

\bibitem{chen1999}
J.~Chen and R.~J.~Patton, \emph{Robust Model-Based Fault Diagnosis for Dynamic
Systems}. Boston, MA: Kluwer Academic Publishers, 1999.

\bibitem{deluca2003fdi}
A.~De~Luca and R.~Mattone, ``Actuator failure detection and isolation using
generalized momenta,'' in \emph{Proc.\ IEEE Int.\ Conf.\ Robotics and
Automation (ICRA)}, 2003, pp.~634--639.

\bibitem{visinsky1995fdi}
M.~L.~Visinsky, J.~R.~Cavallaro, and I.~D.~Walker, ``A dynamic fault tolerance
framework for remote robots,'' \emph{IEEE Trans.\ Robotics and Automation},
vol.~11, no.~4, pp.~477--490, 1995.

\bibitem{georgopoulos1986}
A.~P.~Georgopoulos, A.~B.~Schwartz, and R.~E.~Kettner, ``Neuronal population
coding of movement direction,'' \emph{Science}, vol.~233, no.~4771,
pp.~1416--1419, 1986.

\bibitem{kajita2003}
S.~Kajita \emph{et al.}, ``Biped walking pattern generation by using preview
control of zero-moment point,'' in \emph{Proc.\ IEEE Int.\ Conf.\ Robotics and
Automation (ICRA)}, 2003, pp.~1620--1626.

\bibitem{pratt2006capture}
J.~Pratt, J.~Carff, S.~Drakunov, and A.~Goswami, ``Capture point: a step toward
humanoid push recovery,'' in \emph{Proc.\ IEEE-RAS Int.\ Conf.\ Humanoid
Robots}, 2006, pp.~200--207.

\bibitem{wieber2006}
P.-B.~Wieber, ``Trajectory free linear model predictive control for stable
walking in the presence of strong perturbations,'' in \emph{Proc.\ IEEE-RAS
Int.\ Conf.\ Humanoid Robots}, 2006, pp.~137--142.

\bibitem{kuindersma2016}
S.~Kuindersma \emph{et al.}, ``Optimization-based locomotion planning, control,
and simulation studies with the Atlas humanoid robot,'' \emph{Autonomous
Robots}, vol.~40, no.~3, pp.~429--455, 2016.

\bibitem{hwangbo2019}
J.~Hwangbo \emph{et al.}, ``Learning agile and dynamic motor skills for legged
robots,'' \emph{Science Robotics}, vol.~4, no.~26, 2019.

\bibitem{kalakrishnan2010}
M.~Kalakrishnan, J.~Buchli, P.~Pastor, M.~Mistry, and S.~Schaal, ``Fast,
robust quadruped locomotion over unknown terrain,'' in \emph{Proc.\ IEEE Int.\
Conf.\ Robotics and Automation (ICRA)}, 2010, pp.~4672--4677.

\bibitem{grandia2021}
R.~Grandia, F.~Farshidian, R.~Ranjan, and M.~Hutter, ``Hierarchical MPC for
robust multi-gait locomotion of legged robots,'' \emph{IEEE Robotics and
Automation Letters}, vol.~6, no.~2, pp.~1471--1478, 2021.

\bibitem{fankhauser2018}
P.~Fankhauser \emph{et al.}, ``Robust fast navigation of autonomous mobile
robots in natural environments,'' \emph{Autonomous Robots}, vol.~42, no.~7,
pp.~1435--1452, 2018.

\end{thebibliography}

\end{document}
